% ============================================================================
%  model_simpleswing: specification
%  The standard uniform swing model.
%  Source: scripts/model/live/R/model_swing.R.
%  Compile: pdflatex swing_model_spec.tex
% ============================================================================
\documentclass[11pt]{article}

\usepackage[margin=1in]{geometry}
\usepackage{amsmath, amssymb}
\usepackage{microtype}
\usepackage{enumitem}
\usepackage[colorlinks=true, allcolors=blue]{hyperref}

\setlist[itemize]{itemsep=2pt, topsep=4pt}
\newcommand{\dem}{\mathrm{dem}}
\newcommand{\rep}{\mathrm{rep}}
\newcommand{\oth}{\mathrm{oth}}
\newcommand{\var}[1]{\texttt{#1}}

\title{\texttt{model\_simpleswing}: specification}
\author{Young-Kyung Kim}
\date{\today}

\begin{document}
\maketitle

\noindent
This documents the standard uniform swing model as implemented in
\var{scripts/model/live/R/model\_swing.R}. It is one of the two baselines the precinct
$\times$ vote-mode model is compared against. Sections follow the same order as the main model's specification so
the two can be read side by side.

\section{Outcomes}

The model works on precinct $\times$ vote-mode cells, with
$m \in \{\text{mail}, \text{early}, \text{eday}\}$, and estimates two things per cell:

\begin{itemize}
  \item a \textbf{turnout} correction, on the log scale, relative to a fixed baseline;
  \item a \textbf{two-party share} correction --- the \emph{swing} --- on the raw share
        scale, relative to the precinct's prior lean.
\end{itemize}

The two are estimated by separate, independent regressions. Third-party share is not
estimated at all: it is held fixed at the precinct's prior value
$tp_i = \min(\max(\var{oth\_base}_i, 0), 0.25)$.

\section{The baseline}

Both corrections are expressed relative to a baseline expectation of votes per cell,
$e_{im}$. The precinct's expected total $\var{est\_turnout}_i$ is split across modes using
the precinct's own 2022 mode history, shrunk toward its county's split:
\begin{equation}
  q_{im} = \frac{v_{im} + k\, c_{\,c(i)m}}{\textstyle\sum_{m'} v_{im'} + k},
  \qquad
  p_{im} = \frac{q_{im}}{\sum_{m'} q_{im'}},
  \qquad
  e_{im} = \max\big(\var{est\_turnout}_i \cdot p_{im},\ 0.5\big),
  \label{eq:base}
\end{equation}
with $k = 40$, $v_{im}$ the precinct's 2022 voters in mode $m$, and $c_{cm}$ the county's
mode share. This is the same baseline \texttt{model\_cory} is given, and the same shrinkage
constant the main model uses for its mode priors, so the three models start from a common
point.

\section{Turnout}

Write $c_{im}$ for the ballots counted so far in a cell and $f_{im}$ for the fraction of
the cell that has reported. The turnout outcome is the log ratio of counted to expected,
\begin{equation}
  \ell_{im} = \log \frac{\max(c_{im},\, 0.5)}{e_{im}},
  \label{eq:lr}
\end{equation}
fitted by two-way fixed effects on county and mode,
\begin{equation}
  \ell_{im} = \alpha_{c(i)} + \gamma_m + \varepsilon_{im},
  \label{eq:fe}
\end{equation}
weighted by $e_{im}$ and estimated with \texttt{fixest::feols}. Only \emph{counted-out}
cells enter the fit: those with $f_{im} \geq 1$ and at least one two-party vote.
Partially-counted cells contribute nothing to the turnout estimate.

Cells whose county or mode carries no fitted value fall back to the weighted mean for
that mode, and then to the overall weighted mean.

The predicted cell total floors the fit at what is already counted:
\begin{equation}
  \tau_{im} =
  \begin{cases}
    c_{im} & \text{if the cell is counted out,}\\[2pt]
    \max\!\big(c_{im},\ e_{im}\exp(\hat\ell_{im})\big) & \text{otherwise.}
  \end{cases}
  \label{eq:tot}
\end{equation}

Note what equations~\eqref{eq:fe} and~\eqref{eq:tot} together imply. The fitted correction
$\hat\ell_{im}$ varies only at the county $\times$ mode level, so \emph{all} precinct-level
variation in predicted turnout comes from the baseline $e_{im}$; the fit supplies a shared
multiplicative adjustment and nothing else. The model therefore never makes a
precinct-specific turnout prediction.

\section{Share}

The share outcome is the swing away from the precinct's prior two-party lean
$\var{dem2}_i$, on the raw share scale rather than the logit scale:
\begin{equation}
  y_{im} = \frac{c^{\dem}_{im}}{\max(c^{\dem}_{im} + c^{\rep}_{im},\, 1)} - \var{dem2}_i,
\end{equation}
fitted by the same two-way specification \eqref{eq:fe} and the same weights. Here the fit
set is wider: \emph{any} cell with at least one two-party vote counts, whether or not it
has finished reporting. The predicted share is clamped,
$s_{im} = \min(\max(\var{dem2}_i + \hat y_{im},\, 0.001),\, 0.999)$.

\section{Partial counts}

Counted votes are kept exactly and the remainder is allocated at the predicted share:
\begin{equation}
  r_{im} = (\tau_{im} - c_{im})(1 - tp_i),
  \qquad
  \dem_{im} = c^{\dem}_{im} + r_{im}\, s_{im},
  \qquad
  \rep_{im} = c^{\rep}_{im} + r_{im}(1 - s_{im}).
\end{equation}
A counted-out cell is substituted exactly, bypassing the prediction entirely.

\section{Uncertainty}

The interval comes from $B = 100$ \textbf{county bootstrap} replicates. Each replicate
resamples the reporting counties with replacement and refits both equations with weights
$w_{im} = e_{im} \cdot \#\{c(i)\}$, where $\#\{c\}$ is the county's multiplicity in that
draw. The resulting statewide quantities are then re-centered on the point fit:
\begin{equation}
  \tilde S^{(b)} = S^{(b)} - \overline{S} + \hat S_{\text{point}}.
  \label{eq:recentre}
\end{equation}

Checkpoints at which nothing has been counted are skipped: with no reported cells the
model has nothing to estimate.

\section*{Limitations}
\addcontentsline{toc}{section}{Limitations}

Stated plainly, and without prejudice to the model's purpose --- it is a deliberately
simple baseline and several of these are the price of that simplicity.

\begin{itemize}
  \item \textbf{No ceiling.} Nothing bounds $\tau_{im}$ by the precinct's registrants.
        In practice this is far less dangerous than it sounds: because the fitted
        correction is constant within a county $\times$ mode cell, the model has very
        little variance available to blow up with. It cannot produce the kind of
        precinct-level overshoot an unconstrained per-precinct prediction can.
  \item \textbf{No adding up.} The three modes of a precinct are predicted independently
        and summed; nothing forces $\sum_m \tau_{im}$ to respect a precinct total. The
        per-cell floor in \eqref{eq:tot} is consistent with this, since there is no
        adding-up property for it to break.
  \item \textbf{Turnout and share are independent.} They are two unrelated fits, and
        turnout reaches the share aggregate only through the size of the remainder. Their
        correlation is not propagated. Measured within-mode $\rho(\tau, s)$ varies in sign
        across modes, so this is not a wash in either direction.
  \item \textbf{Asymmetric information use.} The share equation uses partially-counted
        cells; the turnout equation uses only counted-out ones. Early in the night the
        turnout fit can therefore be resting on very few cells.
  \item \textbf{The interval is a spread, not a posterior.} Re-centering
        \eqref{eq:recentre} means the bootstrap cannot express bias, only dispersion, and
        that dispersion reflects between-county resampling alone --- not residual
        precinct variance and not baseline error. Coverage comparisons against a model
        that draws from a posterior are not like-for-like, and should be read with
        interval \emph{width} alongside.
  \item \textbf{Third party is not modeled.} It is a fixed per-precinct constant, so no
        uncertainty in it reaches the aggregate.
\end{itemize}

\end{document}
